\chapter{Diaphragm Contraceptive Device}

\section*{Definition and Overview}
The contraceptive diaphragm is a reusable, dome-shaped silicone cup that a woman inserts into her vagina before sex to cover the cervix and prevent sperm from reaching the uterus. It is used together with spermicide, which is applied to the diaphragm before insertion. The diaphragm is a barrier method of contraception that is hormone-free, discreet, and can be used as needed, but it must be fitted by a health professional, inserted correctly every time, and left in place for at least six hours after sex. This chapter explains how the diaphragm works, how it is used, and its advantages and disadvantages.

\section*{Epidemiology}
The diaphragm is a less commonly used contraceptive method in Australia than the pill, implants, or condoms, but it remains a valid option for women who prefer a non-hormonal, user-controlled method. Its effectiveness depends heavily on consistent and correct use. With typical use, around 12\% of women using the diaphragm become pregnant within a year, while with perfect use the failure rate is lower. Use has declined since the 1980s with the rise of long-acting methods, but the diaphragm is still chosen by some women and couples.

\section*{Aetiology and Pathophysiology}
\begin{itemize}
    \item \textbf{Mechanism:} the diaphragm covers the cervix, forming a physical barrier to sperm
    \item \textbf{Spermicide:} a spermicidal cream or gel is applied to the diaphragm and kills or immobilises sperm
    \item \textbf{Fitting:} the diaphragm must be fitted by a doctor or nurse to the correct size and type
    \item \textbf{Insertion:} the woman inserts it into the vagina before intercourse and checks that it covers the cervix
    \item \textbf{Timing:} it must be left in place for at least six hours after the last act of intercourse
    \item \textbf{Reuse:} with proper care the diaphragm lasts about two years before needing replacement
\end{itemize}

\section*{Clinical Features}
\begin{itemize}
    \item The diaphragm is inserted before sex and removed after the required time
    \item It is not felt by most women or their partners during intercourse
    \item It does not affect hormones or interfere with fertility after removal
    \item It requires planning and a supply of spermicide
    \item It is less effective than hormonal methods and implants, especially with inconsistent use
    \item Some women find insertion difficult or develop cystitis or vaginal irritation
\end{itemize}

\section*{Differential Diagnosis}
\begin{itemize}
    \item Other barrier methods, such as male and female condoms and the cervical cap
    \item Hormonal methods, including the pill, implant, and intrauterine systems
    \item Copper intrauterine devices and other non-hormonal options
    \item Natural family planning methods, which require different skills
\end{itemize}

\section*{When to Seek Medical Care}
A diaphragm must be fitted by a health professional, and the fitting should be reviewed after childbirth, significant weight change, or pelvic surgery, as these can change the fit. Seek advice if you are considering the diaphragm, have difficulty with insertion, or have any side effects.

\textbf{Call 000 (or local emergency services) for:}
\begin{itemize}
    \item Severe pelvic pain, fever, or unusual discharge, which may indicate infection
    \item Symptoms of an allergic reaction, such as swelling or difficulty breathing
    \item A retained diaphragm that cannot be removed
\end{itemize}

\section*{Diagnosis and Investigations}
\begin{itemize}
    \item \textbf{Clinical assessment:} discussion of contraceptive needs, medical history, and preferences
    \item \textbf{Fitting:} sizing of the diaphragm and instruction in insertion and removal
    \item \textbf{Check of fit:} ensuring the diaphragm covers the cervix comfortably
    \item \textbf{Follow-up:} review after changes in weight, childbirth, or if problems occur
\end{itemize}

\section*{Management}
\begin{itemize}
    \item \textbf{Use with spermicide:} always apply spermicide to the diaphragm before insertion
    \item \textbf{Correct insertion:} fold and insert the diaphragm, then check it covers the cervix
    \item \textbf{Leave in place:} keep it in for at least six hours after sex; add spermicide before further intercourse without removing it
    \item \textbf{Removal and care:} wash with mild soap and water, dry, and store in its case
    \item \textbf{Replacement:} replace the diaphragm every two years, or sooner if damaged
    \item \textbf{Combined with condoms:} also protects against sexually transmitted infections
\end{itemize}

\section*{Complications}
\begin{itemize}
    \item Pregnancy from incorrect or inconsistent use
    \item Urinary tract infections, which are more common in some diaphragm users
    \item Vaginal irritation or allergy to spermicide or silicone
    \item Displacement during intercourse, reducing effectiveness
    \item Toxic shock syndrome, a rare but serious complication if the diaphragm is left in place too long
\end{itemize}

\section*{Prognosis and Outcomes}
The diaphragm is a safe and effective method of contraception when used correctly and consistently. Most women who choose it and receive good instruction use it successfully for years. Its main limitations are the need for planning and correct technique, which is why it suits motivated users. There are no long-term health risks, and fertility returns immediately when use is stopped.

\section*{Prevention and Self-Care}
\begin{itemize}
    \item Get the diaphragm fitted and recheck the fit when needed
    \item Practise insertion until it is easy and reliable
    \item Always use spermicide with the diaphragm
    \item Remove and clean the diaphragm after use
    \item Do not leave the diaphragm in place for more than 24 hours
    \item Use condoms as well if you need protection from sexually transmitted infections
    \item See a doctor if you have problems or are considering other methods
\end{itemize}

\section*{Sources and Further Reading}
The following reputable sources provide further information for healthcare students and patients seeking reliable, current advice about the contraceptive diaphragm.

\begin{itemize}
    \item Healthdirect Australia. \textit{Diaphragm and cap contraception.} https://www.healthdirect.gov.au/
    \item Family Planning Australia. \textit{Diaphragm contraception: information.} https://www.fpnsw.org.au/
    \item NHS. \textit{Diaphragms and caps: how they work.} https://www.nhs.uk/conditions/contraception/
    \item Mayo Clinic. \textit{Diaphragm: contraception overview.} https://www.mayoclinic.org/
    \item MedlinePlus. \textit{Diaphragm for birth control.} https://medlineplus.gov/
    \item The Royal Women's Hospital. \textit{Contraception choices.} https://www.thewomens.org.au/
\end{itemize}
